\section{The quarter residues for the two original quadratic-eight families}
\label{sec:quadratic8-quartic-residue-strict}

\paragraph{Inputs and scope.}
The class conclusions of AC-G158-class are the inputs to this
strict argument.  Its exact unit arrays are independently checked
below by rational univariate substitution and Horner evaluation.
The original finite characters are the two quadratic characters
of conductor eight times a nontrivial cubic character.
The quartic characters at $269$ occur only in the residual
Kummer splitting fields, and the trivial cubic label is not included.

Put $p=269$, $K=\mathbb Q(\omega)$ and $t=\sqrt[3]p>0$.
Let $\chi:G_K\to\mu_3$ describe $K(t)/K$, and set
$\delta_+=\chi_8|_{G_K}$, $\delta_-=\chi_{-8}|_{G_K}$.
The same symbols denote their specified rational quadratic
extensions when discussing the quartic character fields.
Write $c_p$ for the actual cyclotomic character and $\omega_p$
for its Teichmuller part.  Fix $i\in\mathbb Z_p$ with
$i^2=-1$ and $\bar i=82$.
For any discrete coefficient module $D$ use the full strict group
\[
 \begin{gathered}
 S_{\rm str}(K,D)=\ker\bigl(H^1_{\rm cts}(G_K,D)
                                  \longrightarrow\mathcal L(D)\bigr),\\
 \mathcal L(D)=\prod_{v\mid p}H^1(G_{K_v},D)
              \times\prod_{v\nmid p,\infty}H^1(I_v,D).
 \end{gathered}
 \tag{Q8Q.1}
\]

\begin{theorem}[AC-G158: the quadratic-eight quarter-residue strict rows]
\label{AC-G158}
Let $E/\mathbb Q_p$ be finite and let $O$ be its original ring
of integers, containing the values of $\chi$.
Let $j=1,2$, $s\in\{+,-\}$, $m\in\mathbb Z$, $e\geq1$, and let
$\kappa:G_K\to1+\mathfrak m_O$ be continuous and unramified
at every finite place away from $p$.
If $e+m\equiv66$ or $200\pmod{268}$, then
\[
 \begin{gathered}
 D=(E/O)(\chi^{-j}\delta_s\omega_p^{-m}\kappa^{-1}c_p^{-e}),\\
                          S_{\rm str}(K,D)=0.
 \end{gathered}
 \tag{Q8Q.2}
\]
The corresponding residual strict groups vanish over every finite
field containing their values.  No condition on the absolute
ramification index of the original $O$ is imposed.
For $m=0,\kappa=1$, these are additional weight progressions
for the four already existing finite labels $\chi^{-j}\delta_s$.
\end{theorem}

\begin{proof}
\proofstep{1. The two exact quartic character fields.}
Let $\psi=\omega_p^{67}$, with $\psi(2)=-i$.
For a numerical sign $s=1,-1$, corresponding to $+,-$, put
$\xi_s=\delta_s\psi$ as rational Dirichlet--Galois characters.
They have conductor $8p$; $\xi_+$ is odd and $\xi_-$ is even.
Both squares are $\tau_p=\omega_p^{134}$.
Their cyclic quartic fields, in the actual source orientation, are
\[
 \begin{gathered}
 F_s=\mathbb Q(w_s),\qquad w_s^4+2spw_s^2+100p=0,\\
 L_s=F_s(t),\qquad M_s=KL_s.
 \end{gathered}
 \tag{Q8Q.3}
\]
The field dictionary can also be read directly from the retained
Gaussian-period field $\mathbb Q(v)$ of $\psi$:
\[
 v^4+pv^2+25p=0,\qquad
 w_+=\sqrt2\,v,\qquad w_-=\sqrt{-2}\,v.
 \tag{Q8Q.4}
\]
The quadratic factors have conductor eight and are disjoint from
the period field, which is ramified only at $p$.
In each $C_2\times C_4$ product, the single subgroup generated
by the pair of the quadratic involution and the quartic square
fixes the indicated product in (Q8Q.4).
Its degree is four by the Eisenstein polynomial (Q8Q.3), so it
is exactly the field of $\delta_s\psi$.
The original period field is imaginary; therefore $F_+$ is
imaginary and $F_-$ is totally real.

Let $b$ fix $K$ and the corresponding quadratic factor in (Q8Q.4),
and act on the period field by the generator with $\psi(b)=-i$.
Then
\[
 b(w_s)=-\frac{s}{130}w_s^3-\frac{219}{65}w_s,
 \qquad b^2(w_s)=-w_s,\qquad \xi_s(b)=-i.
 \tag{Q8Q.5}
\]
The CRT representative $809$ is one modulo eight and two modulo
$p$, and verifies this character orientation in both fields.

The unique quadratic subfield of each $F_s$ is
$\mathbb Q(\sqrt p)$, different from $K$.
Thus $F_s\cap K(t)=\mathbb Q$, and
\[
 \begin{gathered}
 \operatorname{Gal}(M_s/\mathbb Q)=S_3\times C_4,
 \qquad[M_s:\mathbb Q]=24,\quad[L_s:\mathbb Q]=12,\\
 \Delta_s=\operatorname{Gal}(M_s/K)=C_3\times C_4.
 \end{gathered}
 \tag{Q8Q.6}
\]
Let $c(t)=\omega t$, $c(w_s)=w_s$; it commutes with $b$.

Let $a$ be the residual coefficient character and let
$r\equiv e+m\pmod{268}$.
The exact inverse-character table for the Kummer quotient is
\[
 \begin{array}{c|c|c}
 r&\eta=a/\bar c_p&\eta^{-1}(b)\\ \hline
 66&\bar\chi^{-j}\bar\xi_s^{-1}&187=-\bar i\\
 200&\bar\chi^{-j}\bar\xi_s&82=\bar i.
 \end{array}
 \tag{Q8Q.7}
\]
Since $j=1,2$, $M_s$ splits $\eta$ with degree twelve over $K$.
The character $\delta_s$ in (Q8Q.2) remains the original
quadratic character; it is not replaced by $\xi_s$.

\proofstep{2. The exact class sector and full-local Kummer descent.}
The input from AC-G158-class is
\[
                            \operatorname{Cl}(L_s)[p]=0
                    \qquad(s=1,-1).
 \tag{Q8Q.8}
\]
If a reflection $r$ in $S_3$ fixes $L_s$, norm to $L_s$
followed by extension is $1+r$ on $\operatorname{Cl}(M_s)[p]$.
It is zero by (Q8Q.8), so $r$ acts by minus one.
Every conjugate reflection has the same property.
Their products generate $C_3$, which consequently acts trivially.
The nontrivial cubic class sector is zero.
No vanishing of the trivial cubic class sector is used.

Let $k=O/\mathfrak m_O$.
Prime-to-$p$ restriction and inflation, including the exact
local decomposition and inertia quotients, identify the residual
strict group with
\[
 \begin{gathered}
 e_{\eta^{-1}}\ker\left(V_{M_s,S}\otimes_{\mathbb F_p}k
             \longrightarrow\mathcal P_{M_s}\otimes_{\mathbb F_p}k\right),\\
 V_{M_s,S}=\{x\in M_s^\times:p\mid v(x)\text{ for }v\nmid p\}/M_s^{\times p},\\
 \mathcal P_{M_s}=\bigoplus_{v\mid p}M_{s,v}^\times/M_{s,v}^{\times p}.
 \end{gathered}
 \tag{Q8Q.9}
\]
The coefficient over $M_s$ is $k(1)$, with additional quotient
twist $\eta$, so invariants select its inverse in Kummer cohomology.
No adjoining of $\mu_p$ is required.
The full local power condition also makes all $p$-adic valuations
divisible by $p$.
The kernel therefore lies in the all-primes valuation group,
whose exact sequence is
\[
 0\longrightarrow\mathcal O_{M_s}^{\times}/\mathcal O_{M_s}^{\times p}
 \longrightarrow V_{M_s}\longrightarrow\operatorname{Cl}(M_s)[p]
 \longrightarrow0.
 \tag{Q8Q.10}
\]
The inverse characters in (Q8Q.7) are nontrivial on $C_3$,
so the vanished class sector leaves only their unit components.

\proofstep{3. Exact local fields, uniformizers and actions.}
Each polynomial (Q8Q.3) is Eisenstein at $p$.
Its degree-four completion and the cubic completion have coprime
total ramification degrees.  Their compositum has degree twelve.
Adding the unramified quadratic $K_p$ gives degree twenty-four,
the full global degree of $M_s$.
Thus each $M_s$ has a single prime above $p$, with indices
$(e_p,f_p)=(12,2)$.

Set $u_s=t/w_s$ at that prime.
The following identities hold exactly, with the same sign $s$:
\[
 \begin{aligned}
 1000000u_s^{12}+s417488u_s^6+269&=0,\\
 w_s&=\frac{52811u_s^3+s125000u_s^9}{61},\qquad t=u_sw_s,\\
 c(u_s)&=\omega u_s,\\
 b(u_s)&=\frac{46711}{15860}u_s+s\frac{6250}{793}u_s^7.
 \end{aligned}
 \tag{Q8Q.11}
\]
Exact substitution verifies $t^3=p$, (Q8Q.3), (Q8Q.5),
$b(t)=t$ and $b^2(u_s)=-u_s$.
The polynomial divided by $1000000$ is Eisenstein.
The valuations of coefficient times $u_s^n$ are distinct modulo
twelve, so its power basis is the full local integral basis over
$\mathcal O_{K_p}$.
In the truncated residue rings the full actions are
\[
 \begin{aligned}
 b(u_+)&=82u_++130u_+^7\pmod{p,u_+^{12}},\\
 b(u_-)&=82u_-+139u_-^7\pmod{p,u_-^{12}}.
 \end{aligned}
 \tag{Q8Q.12}
\]
The seventh-degree terms will be used in the projectors.

\proofstep{4. Two actual units and exact denominator retention.}
Use source unit four in the five-unit file for $L_+$, and source
unit six in the seven-unit file for $L_-$ from AC-G158-class.
Call them $U_+,U_-$.
Their exact tensor vectors use the basis $t^aw_s^b$, index $4a+b$;
their respective common denominators are
\[
                            52455=195p,\qquad104910=390p.
 \tag{Q8Q.13}
\]
The full integer numerator vectors and complete monic integral
degree-twelve relations, both with constant coefficient one,
are supplied without abbreviation in
\path{selected_quadratic8_quartic_units.json}, SHA256
\begin{quote}\ttfamily\small
8589fab6d952f1fdaddaeb4d022d173c\\
da3ac8b41a5471cbb61e97fbe43e13e1.
\end{quote}
The underlying immutable source unit files have SHA256
\begin{quote}\ttfamily\small
cc09226d6f5a1072782b152ac8161a70\\
147000943a6a75d7642e54d2ea889645,\\[2pt]
d3c36fd4771a4cafca7fd0f730591735\\
fcb77990e7edd7a960de22b40f022601.
\end{quote}
These data specify the two global elements and their unit
relations exactly, not just their local approximations.
The verifier substitutes (Q8Q.11) into each complete rational
tensor expression and checks its monic relation by Horner
evaluation in the corresponding rational degree-twelve algebra.
Monicity proves integrality; the constant coefficient one makes
the inverse an integral polynomial in the unit.
All twelve supplied unit relations are checked, but no complete
fundamental-unit basis is assumed in this strict proof.

Only after exact rational conversion does the verifier reduce
modulo $p$.  All resulting coefficients are $p$-integral;
the original factor $p$ in each denominator is retained through
every cancellation.
The resulting twelve coefficients are
\[
 \begin{aligned}
 U_+&\equiv(268,239,88,73,144,41,\\
    &\hspace{3em}74,21,260,210,206,131)\pmod{p,u_+^{12}},\\
 U_-&\equiv(268,174,134,21,75,261,\\
    &\hspace{3em}201,267,164,13,8,207)\pmod{p,u_-^{12}}.
 \end{aligned}
 \tag{Q8Q.14}
\]

\proofstep{5. Four nonzero global-projector local images.}
After removing the Teichmuller residue, the logarithms below
degree twelve are
\[
 \begin{aligned}
 \ell_+&=30u_++16u_+^5+149u_+^7+268u_+^{11},\\
 \ell_-&=95u_-+61u_-^2+90u_-^3+183u_-^4+223u_-^5\\
       &\quad+74u_-^6+56u_-^7+14u_-^8\\
       &\quad+54u_-^9+108u_-^{10}+2u_-^{11}.
 \end{aligned}
 \tag{Q8Q.15}
\]
These are computed by solving $U_s\ell_s'=U_s'$ and integrating
with denominators one through eleven, all invertible modulo $p$.
Apply the actual global correction
$V_s=U_s^3/(U_sc(U_s)c^2(U_s))$.
The operator $3-(1+c+c^2)$ kills the degrees divisible by three
and multiplies the others by three, giving
\[
 \begin{aligned}
 L_+&=90u_++48u_+^5+178u_+^7+266u_+^{11},\\
 L_-&=16u_-+183u_-^2+11u_-^4+131u_-^5\\
    &\quad+168u_-^7+42u_-^8+55u_-^{10}+6u_-^{11}.
 \end{aligned}
 \tag{Q8Q.16}
\]
For $\lambda=82,187$ let
$P_\lambda=\sum_{a=0}^3\lambda^{-a}b^a$ on unit classes modulo $p$.
This is four times the corresponding idempotent.
Using each full nonlinear action from (Q8Q.12) gives
\[
 \begin{aligned}
 P_{82}L_+&=91u_++192u_+^5+246u_+^7+118u_+^{11},\\
 P_{187}L_+&=197u_+^7+139u_+^{11},\\
 P_{82}L_-&=64u_-+255u_-^5+28u_-^7+3u_-^{11},\\
 P_{187}L_-&=106u_-^7+21u_-^{11}.
 \end{aligned}
 \tag{Q8Q.17}
\]
All four values are nonzero, and the verifier checks their
eigenvalue identities under $b$.
Integer representatives of the projector coefficients give
products of global conjugate units before localization.
Changing such representatives by multiples of $p$ does not
change their Kummer classes or the truncated logarithms.

Since $v_{u_s}(p)=12<p-1$, logarithm and exponential identify
local principal units with the maximal ideal.
The logarithm of a $p$th power principal unit has valuation at
least thirteen.  The degree-one and degree-seven terms in
(Q8Q.17) lie below this threshold.
Teichmuller units are themselves $p$th powers, and the root of
a principal unit must have residue one because Frobenius is
bijective on the finite residue field.
Thus all four images are nonzero for the full local Kummer
condition, not merely for a smaller subgroup of units.

\proofstep{6. Unit multiplicities and residual strict zero.}
The twelve archimedean places of $M_s$ above a fixed complex
embedding of $K$ form a regular $\Delta_s$-set.
Its equivariant Dirichlet unit character is
$\mathbb Q[\Delta_s]-\mathbb Q$.
The field cannot contain $\mu_p$, since $268\nmid24$.
As $p\nmid12$, these character multiplicities persist on the
unit quotient modulo $p$.
For each $b$-eigenvalue $82,187$, the nontrivial cubic part has
dimension two over $\mathbb F_p$ and is simple as an
$\mathbb F_p[C_3]$-module, because $X^2+X+1$ is irreducible
over $\mathbb F_{269}$.

Localization is $\Delta_s$-equivariant.
For each sign and eigenvalue, the corresponding nonzero value
in (Q8Q.17) makes localization injective on that simple component.
After scalar extension to $k$ it is injective on both cubic
eigenlines in the selected inverse-character row of (Q8Q.7).
The vanished class sector and (Q8Q.9)--(Q8Q.10) now prove
the residual strict zero for every $j=1,2$ and both signs.

\proofstep{7. Lift over the original coefficient ring.}
Let $\varpi$ be a uniformizer of the original $O$.
The module $D$ is divisible and $D[\varpi]\simeq k(a)$.
At $p$ the quadratic character $\delta_s$ is unramified,
$\kappa$ reduces to one, and residual cubic inertia has order
three.  A residual cyclotomic character has order dividing $268$,
so it cannot cancel that cubic character.
Some inertia element acts by a scalar whose difference from
one is a unit of $O$, and $D^{G_{K_p}}=0$.
Away from $p$, the cyclotomic, Teichmuller and $\kappa$ factors
are unramified.  The remaining finite inertia character has
order dividing six, prime to $p$; its invariants are zero or
the whole divisible module $D$.

For every local group $J$ in (Q8Q.1),
$D^J/\varpi D^J=0$.
The multiplication-by-$\varpi$ long exact sequence gives local
injectivity $H^1(J,D[\varpi])\to H^1(J,D)$.
A global strict class killed by $\varpi$ lifts to a residual
class, and local injectivity makes the lift strict.
It vanishes by Step6.
Every continuous cocycle from compact $G_K$ to discrete $E/O$
has finite image and is killed by a power of $\varpi$.
A nonzero strict class would therefore yield a nonzero one
killed by $\varpi$, a contradiction.
This proves (Q8Q.2) over the original, possibly ramified ring.
\end{proof}

\paragraph{Actual finite labels and the exact additional count.}
For $m=0,\kappa=1$ the four finite labels are
$\chi^{-j}\chi_8$ and $\chi^{-j}\chi_{-8}$, $j=1,2$.
They already occur on the prescribed hosts containing $\zeta_8$
and $K(t)$; no conductor-$269$ quartic character is added.
The new weights are $e=66+134n$, $n\geq0$, so the number of
pairs with $0<e<T$, for positive integer $T$, is
\[
                           4\left\lfloor\frac{T+67}{134}\right\rfloor.
 \tag{Q8Q.18}
\]
Together with the already proved half-residue rows for the same
four labels, their full proved progression is $e\equiv-1\pmod{67}$,
with count $4\lfloor T/67\rfloor$.
The strict zeros do not imply that their individual positive
derived indices vanish.  The local comparison costs and exact
common-window transport are separate arguments.

\paragraph{Evidence and geometric scope.}
The complete finite verifier is
\path{verify_quadratic8_quartic_local_units.py}; its output is
\path{quadratic8_quartic_local_unit_certificate.json}.
The earlier exact field-only calculation is preserved separately.
All source unit arrays, monic relations and selected complete
data are retained in the accompanying files, with exact hashes.
The class theorem is used only for (Q8Q.8) and its actual field
dictionary.  No new analytic replay, full normal-closure class
claim or complete fundamental-unit basis is inferred from the
local computation.
The retained finite-cover and rank-one Artin--Tate dictionary
is SGA~1, Expos\'e~V, \S\S4 and 7; Kummer theory retains the
previously audited Stacks tag~03PK scope.
Class norm/extension and equivariant Dirichlet units retain
their earlier checked scopes.
No trivial cubic label, automorphic endpoint, mixed realization,
all-prime upper bound, H04.4 conclusion or Lean execution is claimed.
